% !TEX root = ../HDR.tex
%%%%%%%%%%%%%%%%%%%%%%%%%%%%%%%%%%%%%%%%%%%%%%%%%%%%%%%%%%%%%%%%%%%%%%%%%%%%%%%%%%
\chapter*{}\addcontentsline{toc}{part}{\IntroductionName}
\vspace{-1.5cm}
    \begin{minipage}[r]{.95\textwidth}\raggedleft
    \HUGE\bfseries\IntroductionName
    \end{minipage}
\vspace{2.5cm}

                    %%%%%%%%%%%%%%%%%%%%%%%%%%%%                  (EDIT BELOW)

\noindent
J'ai écrit ce manuscrit d'Habilitation à Diriger les Recherches comme un livre à destination d'un public de scientifiques intéressés par l'étude de la cognition. 
Il s'adresse à un lecteur informé des sciences cognitives et des techniques informatiques.
Il est rédigé en français mais contient des citations originales en anglais et des termes techniques non traduits, supposant un lecteur bilingue.  

Ma question initiale est celle de la cognition et, pour l'aborder, celle de la conception d'un agent artificiel cognitif. 
Cette question relève du domaine de l'informatique et de l'Intelligence Artificielle (IA), mais elle s'appuie sur une compréhension de la cognition établie par d'autres disciplines.

Plus précisément, cette question se positionne dans le champ d'une IA qui a été qualifiée de \emph{bio-inspirée}, mais les qualificatifs moins utilisés de \emph{philosophie-inspirée} ou \emph{épistémologie-inspirée} pourraient aussi s'appliquer. 
Cette IA se distingue de l'IA basée sur de grandes quantités de données d'entrainement qui occupe le devant de la scène depuis la popularisation des grands modèles de langage.

Les deux premiers chapitres organisent un cheminement théorique et technique qui pose les bases de ma recherche.  
Le chapitre 1 présente les sources d'inspiration philosophiques, épistémologiques et cognitives qui orientent vers le problème de l'apprentissage développemental artificiel précoce.
La revue de l'état de l'art informatique du chapitre 2 inscrit ce problème dans le cadre plus précis de la conception de mécanismes d'organisation des comportements autonomes permettant l'ancrage de la signification dans l'expérience.

Les chapitres suivants présentent ma contribution. 
Le chapitre 3 présente la théorie originale des Événements d'Interaction Valués comme un nouvel angle d'approche permettant de concevoir un agent autonome développemental.
Le chapitre 4 présente des recherches en cours illustrant une méthodologie de conception et d'expérimentation qui s'inscrit dans une logique de rétro-ingénierie de la cognition. 




